%%
%%
%%
\documentclass[a4paper,12pt]{article}

\usepackage{a4wide}
\usepackage{wasysym}
\usepackage{times}
\usepackage{framed}
\usepackage{graphics}
\usepackage{graphicx}

\title{Esercitazione di Embedded Systems}

\begin{document}

\maketitle
\def\X{\textbf{``X''}}
\def\Y{\textbf{``Y''}}
\def\Z{\textbf{``Z''}}
\def\T{\textbf{``T''}}

%%
%%\vspace{1em}
%%

Realizzare un programma che implementi le funzionalit\`a di un
cronometro con la risoluzione di $1/100~secs$.

Il sistema deve operare secondo le specifiche:

\begin{itemize}
\item tasto \X{}, avvia/ferma il cronometro;
\item tasto \Y{}, azzera il cronometro;
\item tasto \Z{}, memory function (vedi oltre);
\item tasto \T{}, visualizza il contenuto della memoria.
\end{itemize}

Il valore del cronometro deve essere visualizzato nel display usando
il formato \textbf{\texttt{SS:CC}}, dove $SS$ sono i secondi e $CC$
sono i $1/100$ di secondo.

Quando il cronometro \`e avviato, i due-punti ($:$) presenti nel
display devono lampeggiare con un periodo di $500~ms$.

\vspace{1em}

La memoria deve funzionare nel seguente modo:
\begin{itemize}
\item Quando il cronometro \`e avviato, la pressione del tasto \Z{}
  implica lo ``snapshot'' del tempo corrente in un array di memoria
  che pu\`o ospitare fino a 10 ``snapshots'';
\item L'operazione di store deve essere segnalata accendendo il LED
  giallo per 1 secondo;
\item Quando il cronometro \`e fermo, \`e possibile visualizzare la
  memoria premendo il tasto \T{}; in tal caso, il sistema accende il
  LED rosso e mostra, sul display, il valore di ogni snapshot per 1
  secondo; alla fine della sequenza il LED rosso si spegne e si
  ritorna al funzionamento normale.
\end{itemize}

%%
\end{document}
%%
